\documentclass[12pt,a4paper]{article}

\usepackage[T2A]{fontenc}
\usepackage[utf8]{inputenc}
\usepackage[russian]{babel}

\usepackage{amsmath,amssymb,amsthm}
\usepackage{geometry}
\usepackage{hyperref}
\usepackage{microtype}
\usepackage{setspace}
\usepackage{caption}
\usepackage{graphicx}

\geometry{
  left=25mm,
  right=25mm,
  top=25mm,
  bottom=25mm
}

\hypersetup{
  colorlinks=true,
  linkcolor=black,
  urlcolor=blue
}

\begin{document}
\selectlanguage{russian}

% ================= ТИТУЛЬНЫЙ ЛИСТ =================

\begin{titlepage}
  \begin{center}
    {Санкт-Петербургский политехнический университет Петра Великого\\
    Институт прикладной математики и механики\\[0.5em]
    \textbf{Высшая школа прикладной математики и вычислительной физики}}

    \vfill

    {\LARGE\bfseries Отчёт по курсовой работе\\ Часть №1\par}
    \vspace{1em}
    {\Large по дисциплине\\[0.3em]
    {\bfseries «Проект по математической статистике»}}

    \vfill

    \begin{flushright}
        Выполнил\\
        студент гр. 5030102/20202:\\
        Зайдиев Артур
    \end{flushright}

    \begin{flushright}
        Проверил\\
        Преподаватель: Заяц Олег Иванович
    \end{flushright}

    \vfill
    Санкт-Петербург\\
    2026
  \end{center}
\end{titlepage}

\setcounter{page}{2}
\newpage

% ================= ОСНОВНАЯ ЧАСТЬ =================

\section*{Вычисление ОИСКО: Показательное распределение и колокольное ядро}

Задана плотность показательного распределения:

\[
f_0(x)=
\begin{cases}
\lambda e^{-\lambda x}, & x \ge 0,\\
0, & x<0
\end{cases}
\]

где параметр распределения принимается равным:

\[
\lambda = 1
\]

Задана плотность колокольного (гауссового) ядра:

\[
K(x)=\frac{1}{\sqrt{2\pi}}e^{-x^2/2}
\]

Найдём характеристические функции.

Для показательного распределения:

\[
E_0(z)=\int_0^\infty \lambda e^{-\lambda x}e^{izx}dx
\]

\[
E_0(z)=\lambda \int_0^\infty e^{-(\lambda-iz)x}dx
\]

\[
E_0(z)=\frac{\lambda}{\lambda-iz}
\]

Следовательно:

\[
|E_0(z)|^2=
\frac{\lambda^2}{\lambda^2+z^2}
\]

Для колокольного ядра:

\[
E_K(z)=\int_{-\infty}^{\infty}
\frac{1}{\sqrt{2\pi}}e^{-x^2/2}e^{izx}dx
\]

\[
E_K(z)=e^{-z^2/2}
\]

---

Относительное интегральное среднеквадратическое отклонение вычисляется по формуле:

\[
\bar{\delta}_n(\xi)=
1+\frac{I_1}{n\xi I_4}
+\left(1-\frac1n\right)\frac{I_2(\xi)}{I_4}
-2\frac{I_3(\xi)}{I_4}
\]

Необходимо вычислить интегралы.

% ================= I1 I4 =================

\subsection*{1. Вычисление $I_1$ и $I_4$}

\[
I_1=
\int_{-\infty}^{\infty}|E_K(z)|^2dz
\]

\[
I_1=
\int_{-\infty}^{\infty}e^{-z^2}dz
\]

Известно, что гауссов интеграл равен:

\[
I_1=\sqrt{\pi}
\]

---

Вычислим $I_4$:

\[
I_4=
\int_{-\infty}^{\infty}|E_0(z)|^2dz
\]

\[
I_4=
\int_{-\infty}^{\infty}
\frac{\lambda^2}{\lambda^2+z^2}dz
\]

Используя табличный интеграл:

\[
\int_{-\infty}^{\infty}\frac{dz}{a^2+z^2}=\frac{\pi}{a}
\]

получаем:

\[
I_4=\pi\lambda
\]

% ================= I2 =================

\subsection*{2. Вычисление $I_2(\xi)$}

\[
I_2(\xi)=
\int_{-\infty}^{\infty}
|E_0(z)|^2|E_K(\xi z)|^2dz
\]

Подставляя выражения:

\[
I_2(\xi)=
\int_{-\infty}^{\infty}
\frac{\lambda^2}{\lambda^2+z^2}e^{-\xi^2z^2}dz
\]

Используем известный интеграл:

\[
\int_{-\infty}^{\infty}
\frac{e^{-az^2}}{z^2+b^2}dz=
\frac{\pi}{b}e^{ab^2}\operatorname{erfc}(b\sqrt a)
\]

Тогда:

\[
I_2(\xi)=
\pi\lambda
e^{\lambda^2\xi^2}
\operatorname{erfc}(\lambda\xi)
\]

где

\[
\operatorname{erfc}(x)=
\frac{2}{\sqrt{\pi}}
\int_x^\infty e^{-t^2}dt
\]

---

% ================= I3 =================

\subsection*{3. Вычисление $I_3(\xi)$}

\[
I_3(\xi)=
\int_{-\infty}^{\infty}
|E_0(z)|^2E_K(\xi z)dz
\]

Подставим:

\[
I_3(\xi)=
\int_{-\infty}^{\infty}
\frac{\lambda^2}{\lambda^2+z^2}
e^{-\xi^2z^2/2}dz
\]

Используя тот же табличный интеграл:

\[
I_3(\xi)=
\pi\lambda
e^{\lambda^2\xi^2/2}
\operatorname{erfc}
\left(
\frac{\lambda\xi}{\sqrt2}
\right)
\]

% ================= ИТОГ =================

\subsection*{Итоговая формула ОИСКО}

Подставляя найденные значения:

\[
I_1=\sqrt{\pi}
\]

\[
I_4=\pi\lambda
\]

\[
I_2(\xi)=
\pi\lambda
e^{\lambda^2\xi^2}
\operatorname{erfc}(\lambda\xi)
\]

\[
I_3(\xi)=
\pi\lambda
e^{\lambda^2\xi^2/2}
\operatorname{erfc}
\left(
\frac{\lambda\xi}{\sqrt2}
\right)
\]

получаем окончательное выражение:

\[
\bar{\delta}_n(\xi)=
1+
\frac{\sqrt{\pi}}{n\xi\pi\lambda}
+
\left(1-\frac1n\right)
e^{\lambda^2\xi^2}
\operatorname{erfc}(\lambda\xi)
\]

\[
-2
e^{\lambda^2\xi^2/2}
\operatorname{erfc}
\left(
\frac{\lambda\xi}{\sqrt2}
\right)
\]

При вычислениях используется:

\[
\lambda=1
\]

% ================= ГРАФИКИ =================

\section*{Построение графиков}

Для анализа поведения ОИСКО построены следующие графики.

\newpage

\subsection*{График 1. Зависимость ОИСКО от параметра сглаживания}

\begin{figure}[ht!]
    \centering
    \includegraphics[width=0.85\textwidth]{figures/f1.png}
    \caption{ОИСКО для различных объёмов выборки}
\end{figure}

\subsection*{График 2. Зависимость оптимального параметра $\xi_{opt}$ от объёма выборки}

\begin{figure}[ht!]
    \centering
    \includegraphics[width=0.85\textwidth]{figures/f2.png}
    \caption{Зависимость $\xi_{opt}(n)$}
\end{figure}

\newpage

\subsection*{График 3. Минимальное значение ОИСКО}

\begin{figure}[ht!]
    \centering
    \includegraphics[width=0.85\textwidth]{figures/f3.png}
    \caption{Нижняя граница ОИСКО}
\end{figure}

\subsection*{График 4. Эффективность ядерной оценки}

\begin{figure}[ht!]
    \centering
    \includegraphics[width=0.85\textwidth]{figures/f4.png}
    \caption{Эффективность оценки}
\end{figure}

\end{document}